% =====================================================================
% TIỂU LUẬN DẪN NHẬP CỦA BẢN CHUYỂN NGỮ
% Không thuộc nguyên tác NCTM.
% =====================================================================

\cleardoublepage
\gdef\CurrentRunningHead{Tiểu luận dẫn nhập}

\chapter*{Tiểu luận dẫn nhập của bản chuyển ngữ}
\addcontentsline{toc}{chapter}{Tiểu luận dẫn nhập của bản chuyển ngữ}

\begin{center}
{\Large\bfseries
Vị trí của \emph{Focus in High School Mathematics: Reasoning and Sense Making}\\
trong giáo dục toán học đương đại}

\vspace{0.65\baselineskip}

{\large
Đọc từ bối cảnh Việt Nam và kỉ nguyên trí tuệ nhân tạo}
\end{center}

\vspace{0.8\baselineskip}

\noindent
\textit{Ghi chú về tư cách văn bản.}
Tiểu luận này do người chuyển ngữ biên soạn và không thuộc nguyên tác của
Hội đồng Quốc gia Giáo viên Toán học Hoa Kì (National Council of Teachers of
Mathematics, NCTM). Những nhận định về bối cảnh Việt Nam, về trí tuệ nhân tạo
và về ý nghĩa hiện thời của cuốn sách là kết quả của việc đặt văn bản năm 2009
vào những hoàn cảnh xuất hiện sau đó; chúng không được quy cho NCTM nếu nguyên
tác không phát biểu như vậy.

\vspace{0.6\baselineskip}

\noindent
\textit{Tóm tắt.}
\emph{Focus in High School Mathematics: Reasoning and Sense Making} (NCTM,
2009) không phải là một chương trình toán trung học thay thế, một danh mục nội
dung hay một sách giáo khoa. Nó là một khung định hướng nhằm đặt suy lí và
kiến tạo ý nghĩa vào nền tảng của việc học toán, qua đó tác động đồng thời đến
cách lựa chọn nội dung, nhiệm vụ, dạy học, đánh giá, cơ hội học tập và sự phối
hợp giữa các thành phần của hệ thống giáo dục. Tiểu luận này xác định vị trí
của FHSM trong dòng phát triển tư tưởng của NCTM; phân tích mức độ mà cuốn
sách có thể được xem là mang một triết lí giáo dục toán học; đối chiếu tầm nhìn
ấy với lịch sử và tiến trình cải cách giáo dục toán học ở Việt Nam; và xem xét
lại giá trị của suy lí, kiến tạo ý nghĩa, chứng minh, diễn giải và đánh giá tính
hợp lí trong bối cảnh trí tuệ nhân tạo tạo sinh. Luận điểm chính của bài viết
là giá trị bền vững của FHSM không nằm ở một tập hợp kĩ thuật dạy học riêng
lẻ, mà ở việc nó thay đổi tiêu chuẩn dùng để nhận biết một trải nghiệm học
toán có chất lượng: từ sự hiện diện của nội dung và việc hoàn thành thủ tục
sang chất lượng của hoạt động trí tuệ mà người học thực sự tham gia. Trong
bối cảnh Việt Nam đang chuyển từ chương trình thiên về mô tả nội dung sang
chương trình định hướng phát triển phẩm chất và năng lực, đồng thời bước vào
môi trường học tập có trí tuệ nhân tạo, khung nhìn này có ý nghĩa tham khảo
đặc biệt; tuy nhiên, nó cần được tiếp nhận như một nguồn tư tưởng để đối thoại
và kiểm định, không như một mô hình ngoại nhập phải sao chép.

\medskip

\noindent
\textit{Từ khóa:} giáo dục toán học; suy lí; kiến tạo ý nghĩa; NCTM; chương
trình giáo dục phổ thông; giáo dục toán học Việt Nam; trí tuệ nhân tạo; đánh
giá; công bằng giáo dục.

\section*{1. Đặt vấn đề: xác định đúng loại hình và tầm vóc của FHSM}

Một cuốn sách giáo dục có thể tạo ảnh hưởng theo nhiều cách. Có cuốn thay đổi
những nội dung được đưa vào chương trình; có cuốn cung cấp phương pháp dạy một
chủ đề cụ thể; có cuốn tổng kết một lĩnh vực nghiên cứu. \emph{Focus in High
School Mathematics: Reasoning and Sense Making}, xuất bản năm 2009, thuộc một
loại khác. Chính trong phần Dẫn nhập, NCTM phân biệt ấn phẩm này với những tài
liệu đương thời chủ yếu phân tích các chủ đề nên được dạy trong từng khóa toán
trung học. FHSM chủ ý chuyển trọng tâm sang những nhấn mạnh về chương trình và
dạy học có khả năng làm cho \emph{reasoning} và \emph{sense making} trở thành
nền tảng của nội dung được dạy và học. NCTM còn mô tả chức năng của ấn phẩm
như một ``bộ lọc then chốt'' để xem xét một cách tổ chức chương trình
(NCTM 2009).

Vì vậy, sẽ không chính xác nếu xem FHSM như một chương trình thay thế cho
chương trình toán trung học, càng không nên đọc nó như một tuyển tập kĩ thuật
sư phạm. Câu hỏi mà nó đặt ra nằm ở một tầng khác: bất kể một hệ thống chọn
dạy những chủ đề nào, học sinh đang được yêu cầu \emph{làm gì về mặt trí tuệ}
khi học những chủ đề ấy? Các em có nhận ra quan hệ và cấu trúc, hình thành
phỏng đoán, tìm căn cứ, kiểm tra tính hợp lí, giải thích và biện minh, đánh giá
lập luận, diễn giải kết quả trong bối cảnh, kết nối tri thức mới với tri thức
đã có hay không? Hay hoạt động học tập chủ yếu dừng ở việc tiếp nhận thủ tục,
nhận dạng dạng bài và tái tạo thao tác trong những tình huống quen thuộc?

Cách đặt vấn đề này giải thích tầm vóc thật sự của cuốn sách. Tầm vóc ấy
không nên được đo bằng độ cao của kiến thức toán học trong các ví dụ. Phần lớn
toán học được sử dụng trong FHSM vẫn nằm trong phạm vi phổ thông. Nó cũng
không phải một chuyên luận triết học xây dựng hệ thống phạm trù từ đầu. Điểm
quan trọng nằm ở \emph{cấp độ can thiệp}: FHSM đề nghị thay đổi tiêu chuẩn
dùng để tổ chức và đánh giá toàn bộ trải nghiệm toán học của học sinh. Khi
suy lí và kiến tạo ý nghĩa được coi là nền tảng thay vì phần bổ sung, sự thay
đổi phải lan sang cách thiết kế nhiệm vụ, cách tổ chức đối thoại trong lớp,
cách sử dụng công nghệ, cách hiểu về chứng minh, cách đánh giá, cách phân phối
cơ hội học tập và cách phối hợp giữa chương trình, dạy học và khảo thí.

Cấu trúc ba phần của FHSM phản ánh đúng tham vọng ấy. Phần thứ nhất xác định
suy lí và kiến tạo ý nghĩa, mô tả các thói quen suy lí và lộ trình phát triển
của chúng. Phần thứ hai đưa hai quá trình này vào năm vùng nội dung lớn của
toán trung học: số và đo lường, kí hiệu đại số, hàm số, hình học, thống kê và
xác suất. Phần thứ ba rời khỏi phạm vi một bài học đơn lẻ để xem xét công bằng,
tính nhất quán giữa chương trình--dạy học--đánh giá, và trách nhiệm của học
sinh, gia đình, giáo viên, nhà quản lí, nhà hoạch định chính sách, giáo dục đại
học và những người phát triển chương trình. Một văn bản bắt đầu từ hoạt động
nhận thức của học sinh rồi kết thúc ở cấu trúc của cả hệ thống giáo dục đã tự
cho thấy rằng đối tượng của nó không chỉ là ``phương pháp dạy học'', mà là
điều kiện để một nền giáo dục toán học tạo ra một loại hiểu biết nhất định.

\section*{2. FHSM trong dòng phát triển tư tưởng của NCTM}

FHSM không xuất hiện như một sáng kiến biệt lập. Dẫn nhập của chính cuốn sách
đặt nó vào một tiến trình dài trong các văn kiện NCTM. \emph{An Agenda for
Action} (1980) đưa giải quyết vấn đề vào vị trí nổi bật. \emph{Curriculum and
Evaluation Standards for School Mathematics} (1989) nhấn mạnh một lõi toán
học chung cùng các quá trình giải quyết vấn đề, suy lí, kết nối và giao tiếp.
\emph{Principles and Standards for School Mathematics} (2000) tổ chức chương
trình bằng năm Chuẩn nội dung và năm Chuẩn quá trình. \emph{Curriculum Focal
Points} (2006) tiếp tục tìm kiếm sự tập trung và nhất quán ở các lớp trước
trung học, trong khi vẫn yêu cầu các quá trình toán học phải thấm vào việc học
nội dung. FHSM có thể được hiểu như một bước tập trung hóa tiếp theo: thay vì
xây thêm một danh mục nội dung, nó đặt suy lí và kiến tạo ý nghĩa vào vai trò
nền tảng cho những quá trình và nội dung đã được xác lập trước đó
(NCTM 1980, 1989, 2000, 2006, 2009).

Điểm này quan trọng vì nó ngăn một cách đọc giản lược: ``NCTM chuyển từ dạy
kiến thức sang dạy tư duy''. FHSM không phủ định nội dung toán học và cũng
không đối lập năng lực với tri thức. Chương 1 nói rõ trọng tâm vào suy lí và
kiến tạo ý nghĩa phải được phát triển \emph{trong bối cảnh của nội dung quan
trọng}; học sinh vẫn phải thực hiện chính xác các thủ tục, nhưng đồng thời
phải hiểu vì sao chúng có hiệu lực, biết khi nào có thể dùng chúng và biết
diễn giải kết quả. Quan điểm này gần với cấu trúc ``năng lực toán học'' trong
\emph{Adding It Up}, nơi sự hiểu biết khái niệm, sự thành thạo thủ tục, năng
lực chiến lược, suy lí thích ứng và khuynh hướng tích cực được xem như những
mạch đan xen của một năng lực thống nhất (Kilpatrick, Swafford và Findell
2001).

Bởi thế, FHSM không tham gia cuộc tranh luận bằng một lựa chọn đơn giản giữa
``kiến thức'' và ``năng lực'', hay giữa ``kĩ năng'' và ``hiểu''. Nó đặt vấn đề
về \emph{cấu trúc của tri thức mà học sinh hình thành}. Một thủ tục có thể
được nhớ như chuỗi hành động rời rạc; cũng thủ tục ấy có thể được đặt trong
mạng lưới quan hệ với biểu diễn, khái niệm, điều kiện áp dụng và lí do khiến
nó đúng. Hai học sinh có thể cho cùng một đáp số nhưng sở hữu hai cấu trúc
tri thức rất khác nhau.

Ví dụ về công thức khoảng cách ở Chương 1 thể hiện sự khác biệt này với độ
tương phản đặc biệt rõ. Trong một kịch bản, học sinh cố khôi phục công thức từ
những mảnh kí hiệu đã nhớ, thậm chí lẫn với công thức trung điểm. Trong kịch
bản khác, bài toán được tái dựng từ tình huống hình học: sự thay đổi theo hai
phương vuông góc dẫn tới tam giác vuông, định lí Pythagoras cung cấp quan hệ,
và công thức tổng quát xuất hiện như kết quả của một chuỗi kết nối. Điều đáng
chú ý không phải giáo viên đã ``giải thích hay hơn'', mà là học sinh được tham
gia vào một loại hoạt động khác: các em có thể tái tạo công thức vì công thức
đã được đặt trong một mạng lưới lí do và quan hệ. FHSM dùng ví dụ ấy để cho
thấy suy lí và kiến tạo ý nghĩa vừa là kết quả của dạy học, vừa là phương
tiện để học sinh đi đến hiểu biết về toán học (NCTM 2009).

Từ đây có thể hiểu tại sao Chương 2 nhấn mạnh rằng suy lí không phải một nhóm
chủ đề mới để thêm vào một chương trình vốn đã đông nội dung, mà là một
\emph{lập trường đối với việc học toán}. Một lập trường như vậy không hỏi
trước hết ``đã dành bao nhiêu tiết cho suy lí'', mà hỏi liệu trong mỗi lĩnh
vực nội dung, học sinh có thường xuyên phải tìm cấu trúc, đưa ra và kiểm tra
phỏng đoán, lựa chọn biểu diễn, xây dựng căn cứ, phản tư về lời giải và khái
quát hóa hay không. Nói cách khác, suy lí không phải một ``nội dung thứ sáu''
nằm cạnh đại số, hình học hay thống kê; nó là một phương thức hoạt động xuyên
qua những nội dung ấy.

\section*{3. Một triết lí giáo dục toán học ở mức nào?}

Câu hỏi liệu FHSM có chứa một triết lí giáo dục hay không cần được trả lời với
sự dè dặt về thuật ngữ. Nếu ``triết lí giáo dục'' được hiểu là một hệ thống
hoàn chỉnh về bản chất của tri thức, con người, giá trị, mục đích tối hậu của
giáo dục và trật tự xã hội, FHSM rõ ràng không phải một chuyên luận như thế.
NCTM không xây dựng một bản thể luận về toán học, không khảo cứu có hệ thống
các trường phái nhận thức luận, cũng không trình bày một triết học chính trị
của giáo dục.

Nhưng nếu hiểu triết lí giáo dục ở nghĩa hẹp hơn: một tập hợp những quan niệm
nhất quán về thế nào là biết toán, người học phải giữ vai trò gì trong việc
hình thành tri thức, điều gì có giá trị trong hoạt động toán học, và hệ thống
giáo dục phải được tổ chức ra sao để những giá trị ấy có thể trở thành hiện
thực, thì FHSM có một lập trường triết lí khá rõ.

Trước hết là một quan niệm về \emph{biết toán}. Đối với FHSM, biết toán không
thể rút gọn vào việc sở hữu một tập hợp kết quả và thủ tục. Suy lí được định
nghĩa rộng như quá trình đi tới kết luận dựa trên chứng cứ hoặc giả định đã
nêu; trong toán học, nó trải từ quan sát quy nạp, giải thích và biện minh phi
hình thức, tìm tòi, phỏng đoán và những khởi đầu sai hướng cho tới suy diễn
hình thức và chứng minh. Kiến tạo ý nghĩa là quá trình phát triển sự hiểu biết
bằng cách kết nối tình huống, bối cảnh hay khái niệm với tri thức đã có.
Như vậy, tri thức toán học có chất lượng không chỉ có ``cái gì'' mà còn có
``vì sao'', ``liên hệ với cái gì'', ``dựa vào đâu'', ``có hiệu lực trong điều
kiện nào'' và ``có nghĩa gì trong bối cảnh''.

Tiếp đó là một quan niệm về \emph{tư cách của người học}. FHSM nhiều lần nhấn
mạnh rằng học sinh phải trực tiếp trải nghiệm suy lí, chứ không chỉ quan sát
giáo viên hay sách giáo khoa suy lí thay mình. Điều này có một hệ quả nhận
thức luận quan trọng: học sinh không được xem chỉ như nơi tiếp nhận tri thức
đã hoàn tất; các em phải trở thành chủ thể có khả năng nêu căn cứ, đánh giá
tính hợp lệ của suy lí của chính mình và của người khác, lựa chọn biểu diễn,
kiểm tra phỏng đoán và chịu trách nhiệm trí tuệ đối với kết luận mình chấp
nhận. Có thể gọi đây là việc nuôi dưỡng \emph{tư cách chủ thể nhận thức}: khả
năng không chỉ sở hữu mệnh đề đúng mà còn tham gia có trách nhiệm vào những
thực hành làm cho một mệnh đề được coi là có căn cứ.

FHSM cũng hàm chứa một quan niệm về quan hệ giữa \emph{phi hình thức và hình
thức}. Cuốn sách bác bỏ một phân đôi đơn giản trong đó kiến tạo ý nghĩa chỉ
thuộc phía trực giác, còn suy lí chỉ bắt đầu khi xuất hiện chứng minh hình
thức. Chương 1 đặt hai quá trình trên một phổ đan xen. Chứng minh được xem là
một cách trình bày suy lí hình thức được xây dựng trên nền tảng kiến tạo ý
nghĩa; đồng thời, chứng minh có thể làm lộ cấu trúc, giải thích vì sao một kết
quả phải đúng và giúp người học đánh giá tính hợp lệ của suy lí. Theo cách
nhìn này, hình thức không thay thế ý nghĩa; mức độ hình thức tăng lên phải
được nâng đỡ bởi cấu trúc ý nghĩa ngày càng giàu hơn. Đây là một lập trường
giáo dục quan trọng, bởi nó chống cả hai xu hướng cực đoan: dạy hình thức như
kí hiệu rỗng, và giữ người học mãi ở trực giác không tiến tới chuẩn mực chứng
minh.

Cuối cùng là một quan niệm về \emph{công bằng và thể chế}. FHSM dành hẳn
Chương 9 cho công bằng, nhưng công bằng ở đây không chỉ là quyền tiếp cận cùng
một nội dung. Vấn đề sâu hơn là quyền tiếp cận những hoạt động toán học có giá
trị nhận thức cao. Nếu một nhóm học sinh được làm những nhiệm vụ yêu cầu
phỏng đoán, giải thích, chứng minh và khái quát, trong khi nhóm khác chỉ được
luyện thủ tục vì bị cho là ``chưa đủ khả năng'', thì hệ thống đã phân phối
không đều chính bản chất của việc học toán. Chương 10 tiếp tục mở rộng vấn đề:
chương trình, dạy học và đánh giá phải tạo thành một chỉnh thể nhất quán.
Chương 11 đưa trách nhiệm ấy sang các bên liên quan. Như vậy, tầm nhìn của
FHSM không dừng ở phẩm chất của một giáo viên giỏi; nó hàm ý rằng một thực
hành nhận thức chỉ có thể trở thành chuẩn mực khi cấu trúc thể chế cùng nâng
đỡ nó.

Những tầng trên cho phép nói rằng FHSM chứa một \emph{triết lí giáo dục toán
học thực hành}: một quan niệm chuẩn tắc về điều gì đáng được coi là học toán
có chất lượng, được triển khai từ hoạt động trí tuệ của cá nhân tới cấu trúc
của hệ thống. Tuy nhiên, phạm vi của kết luận này phải được giữ rõ. FHSM được
viết cho toán trung học và cho những vấn đề của giáo dục Hoa Kì ở cuối thập
niên 2000. Nó không phải một lí thuyết đầy đủ về giáo dục toán học từ mầm non
tới nghiên cứu toán học, càng không phải triết lí duy nhất mà một hệ thống
giáo dục phải theo.

\section*{4. Giáo dục toán học Việt Nam: những lớp lịch sử và một chuyển hướng đang diễn ra}

Đặt FHSM vào Việt Nam đòi hỏi trước hết phải tránh một lối kể lịch sử quá
đơn giản, trong đó giáo dục toán học Việt Nam được mô tả như một nền giáo dục
chỉ có học thuộc và luyện thi đang chờ một mô hình phương Tây đến ``hiện đại
hóa''. Lịch sử toán học và giáo dục toán học Việt Nam có nhiều lớp, nhiều
nguồn ảnh hưởng và nhiều đứt nối hơn thế.

Nghiên cứu của Alexei Volkov về toán học và giáo dục toán học trong Việt Nam
truyền thống cho thấy tri thức toán học đã lưu hành trong một không gian văn
hóa chữ Hán với những quá trình tiếp nhận, biến đổi và trao đổi phức tạp; mô
hình đơn tuyến ``trung tâm--ngoại vi'' không đủ để mô tả lịch sử ấy
(Volkov 2009). Các công trình về toán học Việt Nam truyền thống cũng cho thấy
sự hiện diện của những văn bản tính toán, những thực hành liên quan đến hành
chính, đo đạc và thi cử. Vì thế, lịch sử giáo dục toán học Việt Nam không bắt
đầu cùng nhà trường thuộc địa hay với việc nhập toán học châu Âu.

Nhà trường hiện đại đã làm thay đổi sâu sắc hình thức, nội dung và ngôn ngữ
của giáo dục toán học. Những điều kiện lịch sử của thời thuộc địa, sau đó là
các quan hệ học thuật và giáo dục trong thế kỉ XX, khiến giáo dục toán học
Việt Nam chịu nhiều lớp ảnh hưởng từ các truyền thống Pháp và Nga/Liên Xô.
Trần Cường, Nguyễn Hoàng Anh và Nguyễn Thị Hoài Thu (2024), khi xem xét lịch
sử từ phong trào Toán học mới tới giáo dục toán học Việt Nam, nhấn mạnh rằng
các mối liên hệ ấy có cả tác động tích cực lẫn những hệ quả cần được xem xét
phê phán. Điểm đáng lưu ý không phải tìm một ``nguồn gốc duy nhất'' của
chương trình Việt Nam, mà nhận ra rằng mỗi giai đoạn cải cách đều được thực
hiện trên một nền lịch sử đã tích lũy nhiều lớp quan niệm về toán học, chương
trình và dạy học.

Sau thống nhất đất nước, các đợt cải cách chương trình tiếp tục định hình hệ
thống hiện đại. Bosch, Vu-Nhu và Wijayanti (2023), trong một nghiên cứu về sự
chuyển vị tri thức qua các cải cách chương trình, nhắc tới việc năm 1990 hình
thành chương trình giáo dục phổ thông 12 năm thống nhất và bắt buộc trên toàn
quốc; cải cách năm 2006 đưa vào chương trình cơ bản và nâng cao, đồng thời
thay đổi vị trí của một số nội dung như xác suất và thống kê. Phân tích của
các tác giả đặc biệt hữu ích vì nó cho thấy một nguyên tắc rộng hơn: quyết
định ở cấp chương trình không tự động quyết định những hoạt động toán học cụ
thể mà giáo viên và học sinh thực sự tiến hành. Giữa tri thức được công bố là
phải dạy, tri thức được tổ chức trong sách giáo khoa và tri thức thực sự được
dạy luôn có một quá trình chuyển vị.

Điểm này đặc biệt quan trọng khi đọc Chương trình giáo dục phổ thông 2018.
Ở cấp mục tiêu chính thức, chương trình đánh dấu một chuyển hướng rõ sang
phát triển phẩm chất và năng lực. Chương trình tổng thể xác định năng lực toán
học bằng năm thành phần: tư duy và lập luận toán học, mô hình hóa toán học,
giải quyết vấn đề toán học, giao tiếp toán học, và sử dụng công cụ, phương tiện
học toán; đồng thời nhấn mạnh kết nối giữa các ý tưởng toán học, giữa toán học
với thực tiễn và với các môn học khác (Bộ Giáo dục và Đào tạo 2018). Từ năm
2025, cấu trúc kì thi tốt nghiệp trung học phổ thông cũng được Bộ Giáo dục và
Đào tạo công bố theo định hướng đánh giá năng lực, với các dạng câu hỏi mới và
những bối cảnh có ý nghĩa; môn Toán giảm số lượng lệnh hỏi và dành thời gian
cho những dạng thức được thiết kế để đánh giá tư duy ở mức cao hơn
(Cục Quản lí chất lượng 2023).

Vì vậy, nếu đặt FHSM vào Việt Nam hiện nay, câu hỏi không còn là liệu giáo dục
toán học Việt Nam có nên ``chuyển sang năng lực'' hay không. Ở cấp chương
trình chính thức, chuyển hướng ấy đã được xác lập. Vấn đề khó hơn là làm thế
nào những danh từ như ``tư duy'', ``lập luận'', ``mô hình hóa'', ``giải quyết
vấn đề'' và ``giao tiếp'' trở thành những hoạt động có thể quan sát được trong
một bài học cụ thể; làm thế nào nội dung, nhiệm vụ, dạy học và đánh giá cùng
tạo điều kiện cho những hoạt động ấy; và làm thế nào tránh tình trạng mục tiêu
của chương trình thay đổi nhanh hơn cấu trúc tri thức và chuẩn mực hành động
trong lớp học.

Không nên giải quyết câu hỏi này bằng những phán xét tổng quát về giáo viên
Việt Nam. Một nghiên cứu trường hợp của Nguyễn Duyên Thị và Dung Tran về ba
giáo viên toán trung học tham gia nghiên cứu bài học với các nhiệm vụ thách
thức cho thấy niềm tin và tri thức nghề nghiệp của giáo viên có thể thay đổi
khi sự chú ý chuyển mạnh sang toán học của học sinh, những khó khăn và quan
niệm sai của các em, cũng như việc thiết kế và sắp xếp nhiệm vụ
(Nguyen và Tran 2023). Quy mô nghiên cứu nhỏ không cho phép khái quát cho toàn
hệ thống, nhưng chính giới hạn ấy làm nó hữu ích về phương pháp: nó cho thấy
khoảng cách giữa ``đổi mới chương trình'' và ``đổi mới hoạt động học'' phải
được nghiên cứu ở cấp độ của nhiệm vụ, tương tác và tri thức nghề nghiệp,
không thể suy ra chỉ từ văn bản chính sách.

Trong bối cảnh ấy, FHSM có giá trị đối với Việt Nam không phải vì nó cung cấp
một mô hình Hoa Kì để sao chép, mà vì nó đem lại một khung trung gian giữa
mục tiêu chương trình và hoạt động lớp học. Khi chương trình Việt Nam nói tới
``tư duy và lập luận toán học'', FHSM buộc ta hỏi cụ thể hơn: học sinh đang
phân tích điều gì, tìm quan hệ nào, dựa trên căn cứ nào, kiểm tra phỏng đoán
ra sao, phản tư về lời giải thế nào, mức độ hình thức của suy lí phát triển
theo hướng nào? Khi chương trình nói tới giải quyết vấn đề, FHSM đặt việc giải
vào một chu trình bao gồm phân tích, lựa chọn chiến lược, thực hiện có chủ
đích, diễn giải, kiểm tra tính hợp lí, xem lại giả định và khái quát hóa. Khi
chương trình nói tới giao tiếp, FHSM phân biệt việc chỉ báo cáo đáp số với
việc trình bày một lập luận để người khác có thể kiểm tra. Khi chương trình
nói tới sử dụng công cụ, FHSM hỏi công cụ đang mở rộng khả năng khám phá hay
đang thay thế hoạt động nhận thức mà học sinh cần hình thành.

Sự tương hợp này không có nghĩa hai hệ khung là đồng nhất. Năm thành phần năng
lực trong Chương trình 2018 là cấu trúc chính thức của Việt Nam; FHSM không
thay thế chúng. Giá trị thích đáng hơn của FHSM là cung cấp một lăng kính
nhận thức đủ tinh để kiểm tra chất lượng của những hoạt động được thiết kế
nhằm hiện thực hóa các năng lực ấy.

\section*{5. Kỉ nguyên trí tuệ nhân tạo và sự thay đổi quan hệ giữa sản phẩm với năng lực}

FHSM được xuất bản năm 2009. Mọi diễn giải cuốn sách trong bối cảnh trí tuệ
nhân tạo tạo sinh phải bắt đầu bằng sự thật lịch sử ấy. NCTM không thể viết
FHSM như một phản ứng đối với những hệ thống tạo sinh phổ biến từ đầu thập
niên 2020; vì vậy, không nên gán cho văn bản những tiên đoán mà nó chưa từng
đưa ra. Điều có thể làm một cách nghiêm túc là hỏi: những thay đổi công nghệ
hiện nay làm biến đổi điều kiện của giáo dục toán học ra sao, và trong những
điều kiện mới ấy, các luận điểm của FHSM còn giữ giá trị nào?

Sự thay đổi quan trọng nhất không phải máy tính ``biết làm toán'' lần đầu
tiên. Máy tính bỏ túi, hệ đại số máy tính, phần mềm hình học động, công cụ
thống kê và các hệ thống tìm kiếm đã từ lâu tách một phần thao tác ra khỏi
người học. Điểm mới của trí tuệ nhân tạo tạo sinh là phạm vi của sự tách ấy
mở rộng từ tính toán sang những sản phẩm trước đây thường được xem là dấu hiệu
trực tiếp hơn của hoạt động trí tuệ: lời giải bằng ngôn ngữ tự nhiên, giải
thích, lập luận, mã máy tính, mô hình, thậm chí những văn bản mang hình thức
chứng minh. Một sản phẩm có vẻ hoàn chỉnh ngày càng ít cho phép người quan sát
suy ngược một cách chắc chắn rằng người nộp sản phẩm đã thực hiện những hoạt
động nhận thức tương ứng.

Điều này làm nổi lên một phân biệt căn bản giữa \emph{hiệu suất thực hiện
nhiệm vụ} và \emph{sự học}. Một thí nghiệm thực địa của Bastani và cộng sự
(2025) với gần một nghìn học sinh trung học học toán cho thấy quyền truy cập
một công cụ dựa trên GPT-4 có thể làm tăng đáng kể kết quả trong lúc luyện
tập, trong khi nhóm sử dụng giao diện ít ràng buộc lại làm bài kém hơn nhóm
đối chứng khi công cụ bị rút đi ở bài kiểm tra sau đó. Một biến thể gia sư có
các cơ chế bảo vệ việc học, chẳng hạn ưu tiên gợi ý thay vì cung cấp ngay lời
giải, giảm đáng kể tác động bất lợi này. Nghiên cứu ấy không cho phép kết luận
rằng trí tuệ nhân tạo nói chung gây hại cho việc học toán; nó cho thấy một
điều tinh tế hơn và quan trọng hơn: \emph{cùng một năng lực công nghệ có thể
làm tăng thành tích tức thời mà không làm tăng, thậm chí có thể làm giảm, sự
hình thành năng lực khi thiết kế tương tác không bảo vệ hoạt động nhận thức
của người học}.

Nhận định này gặp một cách đáng chú ý với lập trường của NCTM sau khi trí tuệ
nhân tạo tạo sinh đã xuất hiện. Trong tuyên bố \emph{Artificial Intelligence
and Mathematics Teaching} năm 2024, NCTM khẳng định công cụ trí tuệ nhân tạo
không thay thế nhu cầu dạy toán và giải quyết vấn đề; khả năng tạo ra câu trả
lời sai hoặc không hợp lí khiến kiến thức nền và trực giác toán học càng cần
thiết để thẩm định đầu ra. NCTM cũng cho rằng sự tồn tại của những công cụ giải
nhanh tạo áp lực để rời khỏi các đánh giá nông, chỉ đo thao tác, và chuyển một
phần trọng tâm từ ``giải'' sang sự kết hợp giữa ``giải'' và ``kiểm chứng''.
Việc kiểm chứng nhiều lời giải khả dĩ có thể đòi hỏi hiểu biết sâu hơn việc
chỉ tạo ra một lời giải (NCTM 2024).

Ở cấp rộng hơn, hướng dẫn của UNESCO về trí tuệ nhân tạo tạo sinh trong giáo
dục đặt việc bảo vệ quyền và năng lực chủ động của con người vào trung tâm của
một cách tiếp cận lấy con người làm gốc; khung năng lực trí tuệ nhân tạo dành
cho học sinh năm 2024 tiếp tục nhấn mạnh phán đoán phê phán đối với lời giải
do trí tuệ nhân tạo tạo ra, trách nhiệm công dân, hiểu biết nền tảng và khả
năng sử dụng công nghệ một cách an toàn, có ý nghĩa
(Miao và Holmes 2023; Miao, Shiohira và Lao 2024).

Từ những phát triển sau năm 2009 ấy, có thể đưa ra một diễn giải đương đại về
FHSM. Trong môi trường mà một công cụ có thể sản xuất nhanh một đáp số hoặc
một lời giải có vẻ hợp thức, trọng tâm giáo dục phải chuyển mạnh hơn sang
\emph{quan hệ nhận thức giữa người học và sản phẩm}: người học có hiểu những
đối tượng đang được nói tới không; có nhận ra các giả định không; có biết vì
sao một bước hợp lệ không; có thể phát hiện phản ví dụ hoặc trường hợp đặc
biệt không; có biết một kết quả thống kê cho phép suy luận tới đâu không; có
thể đối chiếu các biểu diễn, kiểm tra tính hợp lí, sửa một lập luận sai, hay
tái dựng một ý tưởng mà không cần sao chép chuỗi thao tác do công cụ cung cấp
không?

Những câu hỏi này nằm rất gần với các thói quen suy lí của FHSM. Tuy nhiên,
mối liên hệ cần được phát biểu đúng chiều: không phải FHSM ``đã tiên đoán AI'',
mà là AI làm rõ hơn một vấn đề FHSM vốn đã đặt ra về sự khác nhau giữa thực
hiện thủ tục và hiểu cấu trúc làm thủ tục có nghĩa. Nếu trước đây sự khác nhau
ấy có thể bị che bởi việc một học sinh làm đúng nhiều bài quen thuộc, công cụ
tạo sinh khiến phép đồng nhất ``sản phẩm đúng = năng lực đã hình thành'' trở
nên còn kém đáng tin hơn.

Từ đây cũng xuất hiện một hệ quả đối với đánh giá. Cấm mọi công cụ không thể
là chiến lược duy nhất, vì nhiều hoàn cảnh học tập và nghề nghiệp thực tế sẽ
bao gồm công cụ. Nhưng cho phép công cụ mà vẫn giữ nguyên loại nhiệm vụ cũ
cũng không đủ. Cần phân biệt rõ mục đích của từng hoạt động: khi nào phải kiểm
tra năng lực độc lập; khi nào cần đánh giá khả năng sử dụng công cụ; khi nào
cần yêu cầu người học thẩm định đầu ra, so sánh chiến lược, truy tìm lỗi, nêu
điều kiện áp dụng, xây dựng phản ví dụ, hoặc giải thích vì sao một kết quả có
thể được chấp nhận. Nói cách khác, trí tuệ nhân tạo không làm biến mất vấn đề
đánh giá suy lí; nó làm cho việc thiết kế bằng chứng về suy lí trở nên khó và
quan trọng hơn.

\section*{6. Tầm nhìn xã hội của FHSM: công dân, công việc, khoa học và công bằng}

Một cách đọc chỉ tập trung vào phương pháp lớp học sẽ bỏ qua chiều xã hội của
FHSM. Ngay đầu Chương 1, NCTM đặt toán trung học trong ba miền chuẩn bị rộng:
toán học cho đời sống, cho nơi làm việc, và cho cộng đồng khoa học--kĩ thuật.
Từ đó, suy lí và kiến tạo ý nghĩa không chỉ được biện minh bằng việc giúp học
sinh làm bài tốt hơn. NCTM liên hệ chúng với khả năng đưa ra quyết định có căn
cứ trong đời sống cá nhân và công cộng, với khả năng thích ứng trong công việc,
và với việc tiếp tục học ở những lĩnh vực khoa học và kĩ thuật (NCTM 2009).

Ở đây có một quan niệm quan trọng về ``năng lực toán học'' mà kỉ nguyên dữ
liệu và trí tuệ nhân tạo chỉ làm rõ thêm. Công dân hiện đại thường không phải
tự mình thực hiện mọi phép tính; nhưng họ phải đọc kết quả do hệ thống khác
tạo ra, hiểu đại lượng đo điều gì, nhận biết mô hình dựa trên giả định nào,
phân biệt tương quan với quan hệ nhân quả, đánh giá mức bất định và phạm vi
suy luận, nhận ra khi một con số được trình bày với độ chính xác không hợp lí,
và quyết định mức độ tin cậy có thể dành cho một kết luận. Đây đều là những
hoạt động mà FHSM đưa vào khung suy lí, đặc biệt rõ trong chương thống kê và
xác suất.

Chiều xã hội ấy còn thể hiện ở quan niệm về công bằng. Nếu năng lực suy lí là
điều cần thiết cho công dân, nghề nghiệp và học tập tiếp tục, thì không thể
xem suy lí như đặc quyền của học sinh ở các lớp nâng cao. Chương 9 của FHSM
đặt vấn đề về phân luồng, nhân khẩu học, kì vọng, niềm tin và thiên kiến bởi
chúng ảnh hưởng tới loại cơ hội toán học học sinh được nhận. Quan điểm này làm
thay đổi nghĩa của ``công bằng trong giáo dục toán học'': ngoài quyền tiếp cận
trường lớp và nội dung, còn có quyền tiếp cận những nhiệm vụ đòi hỏi tư duy,
những cuộc đối thoại trong đó lí lẽ được coi trọng, và sự hỗ trợ đủ để người
học có thể tiến từ suy lí phi hình thức tới những hình thức chặt chẽ hơn.

Điều này đặc biệt đáng lưu ý khi vận dụng FHSM ở Việt Nam. Một hệ thống có thể
đặt mục tiêu năng lực rất cao trên văn bản nhưng vẫn tạo ra chênh lệch nếu
những trải nghiệm giàu suy lí chỉ xuất hiện ở một số trường, lớp, nhóm học
sinh hoặc dưới dạng bồi dưỡng chọn lọc. Tiếp nhận FHSM một cách nghiêm túc vì
thế không thể chỉ là đưa thêm bài khó. ``Khó'' và ``giàu suy lí'' không đồng
nghĩa. Một bài toán kĩ thuật rất khó có thể chỉ yêu cầu nhận dạng và thao tác
quen thuộc đối với người đã luyện đúng dạng; ngược lại, một bài toán kiến thức
không cao vẫn có thể buộc học sinh tìm cấu trúc, lựa chọn mô hình, biện minh
và đánh giá nhiều cách tiếp cận. Công bằng về suy lí trước hết là công bằng về
\emph{chất lượng hoạt động trí tuệ} mà hệ thống tạo cơ hội cho học sinh tham
gia.

\section*{7. Từ chương trình chính thức tới lớp học: giá trị của FHSM như một khung trung gian}

Một trong những khó khăn lớn nhất của cải cách giáo dục nằm ở khoảng cách
giữa ngôn ngữ ở cấp chính sách và hành động ở cấp lớp học. Các khái niệm như
``năng lực'', ``tư duy'', ``giải quyết vấn đề'', ``mô hình hóa'' hay ``giao
tiếp'' có thể được tuyên bố rõ trong chương trình nhưng vẫn để ngỏ vô số câu
hỏi về nội dung, nhiệm vụ, tiến trình học và chuẩn mực đánh giá. Nghiên cứu về
chuyển vị tri thức trong các cải cách chương trình cho thấy những quyết định
lớn không tự động sinh ra một tổ chức tri thức nhất quán ở cấp hoạt động
(Bosch, Vu-Nhu và Wijayanti 2023).

FHSM có giá trị đặc biệt ở khoảng giữa này. Nó đủ khái quát để không trở thành
một giáo án, nhưng đủ cụ thể để người đọc nhìn thấy một khái niệm như ``suy
lí'' vận hành khác nhau trong năm miền nội dung. Với số và đo lường, suy lí
liên quan đến tính hợp lí của đại lượng, độ chính xác và quan hệ định lượng.
Với kí hiệu đại số, nó đi qua việc dùng kí hiệu có ý thức, nhận ra cấu trúc
và giải dựa trên suy lí thay vì thao tác mù. Với hàm số, nó gắn với quan hệ
đồng biến thiên, nhiều biểu diễn và sự chuyển qua lại giữa chúng. Với hình
học, nó bao gồm trực quan, phỏng đoán, lập luận, nhiều cách tiếp cận và chứng
minh. Với thống kê và xác suất, nó phải đối diện trực tiếp với biến thiên,
ngẫu nhiên, mô hình, bất định, phân phối lấy mẫu và phạm vi suy luận.

Cách tổ chức ấy làm lộ một nguyên tắc quan trọng: ``suy lí toán học'' không
phải một kĩ năng chung hoàn toàn tách khỏi tri thức chuyên biệt. Ta không thể
dạy suy lí bằng những lời khuyên trừu tượng rồi giả định học sinh sẽ tự chuyển
nó sang mọi lĩnh vực. Suy lí luôn cần một đối tượng, một hệ biểu diễn, những
chuẩn mực chứng cứ và những quan hệ đặc thù của miền nội dung. Khả năng đánh
giá một phỏng đoán hình học, diễn giải một tham số thống kê và nhận ra cấu trúc
của một biểu thức đại số đều là suy lí, nhưng chúng không đồng nhất về tri
thức cần huy động.

Đây cũng là lí do FHSM cần được đọc cùng Chương 10 về tính nhất quán. Nếu lớp
học khuyến khích học sinh giải thích nhưng bài kiểm tra chỉ thưởng cho tốc độ
và đáp số; nếu chương trình nói tới mô hình hóa nhưng sách và thời lượng buộc
giáo viên tập trung vào chuỗi dạng bài; nếu công nghệ được khuyến khích trong
dạy học nhưng bị tách khỏi mọi nhiệm vụ đánh giá; thì học sinh nhận những tín
hiệu mâu thuẫn về điều hệ thống thật sự coi trọng. Tầm nhìn FHSM chỉ có ý
nghĩa khi chương trình, dạy học và đánh giá cùng hướng tới một chuẩn nhận thức
tương thích.

Vì vậy, chức năng thích hợp nhất của FHSM trong bối cảnh Việt Nam là một
\emph{khung trung gian để chất vấn việc hiện thực hóa chương trình}. Nó không
thay văn bản Chương trình 2018; nó giúp đặt những câu hỏi chi tiết hơn đối với
cách chương trình ấy đi vào sách, bài học và đánh giá. Chính ở tầng này, một
văn bản nước ngoài có thể hữu ích mà không trở thành mô hình để sao chép.

\section*{8. Ý nghĩa của bản chuyển ngữ: từ tiếp cận văn bản tới hình thành ngôn ngữ chuyên môn}

Một bản dịch của FHSM có giá trị trước hết vì mở khả năng đọc một văn bản
quan trọng bằng tiếng Việt. Nhưng đối với một tác phẩm mà phần lớn sức nặng
nằm trong những phân biệt nhận thức tinh tế, chuyển ngữ không chỉ là thay từ.
Nếu những khái niệm gần nhau trong nguyên tác bị hợp nhất trong tiếng Việt,
người đọc có thể tiếp nhận một hệ tư tưởng khác với hệ mà NCTM thực sự xây
dựng.

Quá trình hiệu chỉnh bản dịch này vì thế đã phải phân biệt có hệ thống
\emph{reasoning}, \emph{inference}, \emph{argument},
\emph{argumentation}, \emph{justification}, \emph{proof} và
\emph{explanation}; đồng thời tách \emph{sense making} khỏi
\emph{understanding}, \emph{meaning} và \emph{interpretation}. Kết quả là
một hệ tiếng Việt trong đó \emph{reasoning} được gọi là ``suy lí'',
\emph{sense making} là ``kiến tạo ý nghĩa'', \emph{inference} là ``suy
luận'', \emph{argument} là ``lập luận'', \emph{justification} là ``biện
minh'', \emph{proof} là ``chứng minh'', và \emph{interpretation} là ``diễn
giải''. Lí do của từng lựa chọn được trình bày riêng trong phần
\emph{Khảo luận về hệ thuật ngữ của bản chuyển ngữ} ở cuối sách; ở đây chỉ
cần nhấn mạnh hệ quả học thuật của công việc ấy.

Ngôn ngữ chuyên môn không phải lớp nhãn đặt lên tư tưởng sau khi tư tưởng đã
hoàn tất. Những phân biệt mà một ngôn ngữ có thể biểu đạt rõ sẽ ảnh hưởng tới
những câu hỏi mà cộng đồng chuyên môn có thể đặt ra một cách chính xác. Nếu
``suy lí'', ``suy luận'' và ``lập luận'' được nhập làm một, rất khó thảo luận
rõ việc học sinh đang thực hiện một hoạt động tư duy rộng, đang đi từ dữ liệu
tới kết luận thống kê, hay đang trình bày một cấu trúc lí lẽ để người khác
đánh giá. Nếu ``kiến tạo ý nghĩa'' và ``sự hiểu biết'' cùng bị rút về một chữ
``hiểu'', ta khó phân biệt quá trình hình thành quan hệ với trạng thái hiểu
biết được phát triển nhờ quá trình ấy.

Vì vậy, bản chuyển ngữ có thể đóng một chức năng khiêm tốn nhưng nền tảng:
tạo điều kiện để những cuộc thảo luận về giáo dục toán học diễn ra bằng tiếng
Việt với độ phân giải khái niệm cao hơn. Chức năng này càng đáng lưu ý khi
Chương trình 2018 đã chính thức đưa ``tư duy và lập luận toán học'' vào cấu
trúc năng lực. Muốn hiện thực hóa một mục tiêu như vậy, cộng đồng giáo dục cần
không chỉ một thuật ngữ trong văn bản chương trình mà còn một ngôn ngữ đủ rõ
để mô tả các mức suy lí, loại căn cứ, hoạt động giải thích, biện minh, chứng
minh, diễn giải và phản tư trong lớp học.

Tuy nhiên, bản dịch cũng phải giữ ranh giới của mình. Nó không được dùng các
quan niệm riêng của người chuyển ngữ để làm cho NCTM nói điều NCTM không nói.
Chính vì vậy, phần bình giải được tách khỏi bản tiếng Việt và những khảo luận
của người chuyển ngữ được ghi rõ tư cách văn bản. Nguyên tắc này đặc biệt
quan trọng ở những vấn đề như trí tuệ nhân tạo hay sự mở rộng từ toán phổ
thông lên toán đại học: chúng là những hướng đối thoại được gợi ra từ FHSM,
không phải những phần đã có sẵn trong lập trường năm 2009 của NCTM.

\section*{9. Những giới hạn cần giữ khi tiếp nhận FHSM}

Đánh giá một văn bản có tầm ảnh hưởng không đòi hỏi biến nó thành một kinh
điển bất khả nghi vấn. Trái lại, việc xác định giới hạn là điều kiện để sử
dụng nó đúng.

Giới hạn thứ nhất là lịch sử. FHSM được viết trong bối cảnh Hoa Kì, gắn với
những tranh luận về chương trình trung học, phân luồng, chuẩn, đánh giá, giáo
dục đại học và nhu cầu lao động ở thời điểm cuối thập niên 2000. Một số vấn
đề có tính phổ quát tương đối; một số khác gắn chặt với cơ cấu trường học Hoa
Kì. Khi đọc ở Việt Nam, cần chuyển câu hỏi chứ không sao chép biểu hiện bề
mặt. Chẳng hạn, phân luồng trong trường trung học Hoa Kì và những cơ chế phân
hóa học sinh ở Việt Nam không thể mặc nhiên coi là cùng một hiện tượng, dù cả
hai đều có thể được khảo sát bằng câu hỏi chung về cơ hội tiếp cận hoạt động
toán học chất lượng cao.

Giới hạn thứ hai là cấp học. FHSM tập trung vào trung học phổ thông. Nó có
những phát biểu về sự liên thông từ các lớp trước tới giáo dục sau trung học,
nhưng không cung cấp một lí thuyết đầy đủ về sự trưởng thành toán học ở đại
học, về trừu tượng hóa cấu trúc, về chuyển tiếp sang chứng minh ở toán cao
cấp, hay về hoạt động đặt vấn đề và kiến tạo lí thuyết trong nghiên cứu. Do
đó, ``suy lí và kiến tạo ý nghĩa'' có thể là nền tảng quan trọng cho một tầm
nhìn rộng hơn về giáo dục toán học, nhưng không nên tự động được tuyên bố là
triết lí đầy đủ cho mọi trình độ.

Giới hạn thứ ba là bằng chứng. FHSM là một văn bản định hướng chính thức của
NCTM, được xây dựng trên nghiên cứu và có Thư mục chú giải, nhưng bản thân nó
không phải một tổng quan hệ thống theo chuẩn phương pháp luận hiện nay và
không phải mọi khuyến nghị đều được kiểm chứng bằng cùng một loại bằng chứng.
Khi chuyển từ các nguyên tắc của FHSM sang quyết định chính sách hoặc thiết
kế dạy học trong một bối cảnh khác, cần tiếp tục đối chiếu với nghiên cứu thực
nghiệm, với đặc điểm của người học và với điều kiện thể chế cụ thể.

Giới hạn thứ tư liên quan trực tiếp đến trí tuệ nhân tạo. Mối liên hệ giữa
FHSM và AI trong tiểu luận này là một \emph{diễn giải đương đại}. Nó được nâng
đỡ bởi những vấn đề mà NCTM 2024, UNESCO 2023--2024 và các nghiên cứu gần đây
đặt ra, nhưng không được dùng để viết lại lịch sử của FHSM. Giá trị của phép
đọc này nằm ở chỗ kiểm tra sức sống của một khung tư tưởng cũ dưới điều kiện
mới, không ở việc biến văn bản năm 2009 thành lời tiên tri.

Những giới hạn ấy không làm giảm giá trị của FHSM. Chúng xác định đúng cách
sử dụng nó: như một văn bản có sức tổ chức tư duy, cung cấp tiêu chuẩn và câu
hỏi để phân tích giáo dục toán học, nhưng luôn cần được đặt trong đối thoại
với chương trình quốc gia, nghiên cứu mới và những biến đổi của xã hội.

\section*{10. Kết luận}

Có thể xác định đóng góp trung tâm của \emph{Focus in High School
Mathematics: Reasoning and Sense Making} bằng một chuyển dịch trong đối tượng
quan sát. Một hệ thống giáo dục thường dễ nhìn thấy những gì có thể kiểm kê:
chủ đề đã dạy, số tiết, số bài tập, đáp số, điểm thi, tỉ lệ hoàn thành. FHSM
buộc sự chú ý chuyển sang một đối tượng khó quan sát hơn nhưng gần với bản
chất của việc học: \emph{cấu trúc của hoạt động trí tuệ mà học sinh thực sự
thực hiện trên nội dung toán học}. Từ đối tượng ấy, các vấn đề về nhiệm vụ,
giải thích, chứng minh, công nghệ, đánh giá, công bằng và tính nhất quán của hệ
thống nối lại với nhau.

Đặt trong lịch sử NCTM, FHSM là sự tập trung hóa một dòng tư tưởng đã phát
triển qua nhiều thập niên: toán học nhà trường không thể chỉ được tổ chức bởi
nội dung; các quá trình toán học phải nằm trong chính cách nội dung được học.
Đặt trong bối cảnh Việt Nam, nó gặp một hệ thống đang chính thức chuyển sang
phát triển năng lực nhưng vẫn phải giải quyết công việc khó nhất của mọi cải
cách: chuyển ngôn ngữ của chương trình thành những hoạt động toán học có chất
lượng trong lớp học và thành những hình thức đánh giá tương thích. Đặt trong
kỉ nguyên trí tuệ nhân tạo, nó gặp một điều kiện mới trong đó sản phẩm có vẻ
đúng ngày càng có thể được tạo ra mà không bảo đảm người sử dụng đã hình thành
năng lực tương ứng. Trong cả ba phép đặt ấy, suy lí và kiến tạo ý nghĩa không
xuất hiện như khẩu hiệu thay cho nội dung, mà như tiêu chuẩn để hỏi nội dung
đã trở thành tri thức của người học bằng cách nào.

Bởi thế, giá trị của FHSM đối với giáo dục toán học Việt Nam không nằm ở việc
nó là một văn bản của một tổ chức có uy tín ở Hoa Kì. Giá trị ấy chỉ xuất hiện
khi cuốn sách giúp chúng ta đặt được những câu hỏi tốt hơn đối với thực hành
của chính mình: một nhiệm vụ đang yêu cầu học sinh nghĩ gì; một lời giải bộc
lộ loại hiểu biết nào; một công cụ đang hỗ trợ hay thay thế hoạt động nhận
thức; một phép đánh giá đang đo sản phẩm hay năng lực; một nhóm học sinh có
được tiếp cận những hoạt động toán học giàu suy lí như nhóm khác hay không;
và các thành phần của hệ thống có đang cùng gửi tới học sinh một thông điệp
nhất quán về điều gì có giá trị trong toán học hay không.

Nếu được đọc theo cách ấy, FHSM không cần được nâng thành một hệ triết học
toàn năng để có tầm vóc. Nó làm một việc cụ thể hơn: cung cấp một khung đủ
chặt để chuyển câu hỏi ``dạy toán gì'' sang đồng thời câu hỏi ``con người
được hình thành như thế nào khi học toán''. Trong một nền giáo dục đang đổi
mới theo hướng năng lực và trong một thời đại mà máy móc ngày càng có thể thực
hiện nhiều phần của công việc tạo ra lời giải, câu hỏi thứ hai không thay thế
câu hỏi thứ nhất; nó trở thành điều kiện để câu hỏi thứ nhất còn giữ được ý
nghĩa giáo dục.

\section*{Tài liệu tham khảo của tiểu luận}

\begingroup
\small
\setlength{\parindent}{0pt}
\setlength{\parskip}{0.55\baselineskip}

Bastani, H., Bastani, O., Sungu, A., Ge, H., Kabakcı, Ö., và Mariman, R.
(2025).
``Generative AI without guardrails can harm learning: Evidence from high
school mathematics.''
\emph{Proceedings of the National Academy of Sciences}, 122(26),
e2422633122.
doi: 10.1073/pnas.2422633122.

Bộ Giáo dục và Đào tạo (2018).
\emph{Chương trình giáo dục phổ thông}, ban hành kèm theo Thông tư
32/2018/TT-BGDĐT ngày 26 tháng 12 năm 2018.

Bosch, M., Vu-Nhu, T.-H., và Wijayanti, D. (2023).
``Curriculum Reforms and the Construction of the Knowledge to Be Taught.''
Trong Y. Shimizu và R. Vithal (chủ biên),
\emph{Mathematics Curriculum Reforms Around the World: The 24th ICMI Study},
tr. 101--111. Springer.
doi: 10.1007/978-3-031-13548-4\_7.

Cục Quản lí chất lượng, Bộ Giáo dục và Đào tạo (2023).
``Cấu trúc định dạng đề thi tốt nghiệp THPT từ năm 2025'', công bố ngày
29 tháng 12 năm 2023.

Kilpatrick, J., Swafford, J., và Findell, B. (chủ biên) (2001).
\emph{Adding It Up: Helping Children Learn Mathematics}.
Washington, DC: National Academy Press.

Miao, F., và Holmes, W. (2023).
\emph{Guidance for Generative AI in Education and Research}.
Paris: UNESCO.

Miao, F., Shiohira, K., và Lao, N. (2024).
\emph{AI Competency Framework for Students}.
Paris: UNESCO.
doi: 10.54675/JKJB9835.

National Council of Teachers of Mathematics (NCTM) (1980).
\emph{An Agenda for Action: Recommendations for School Mathematics of the
1980s}. Reston, VA: NCTM.

National Council of Teachers of Mathematics (NCTM) (1989).
\emph{Curriculum and Evaluation Standards for School Mathematics}.
Reston, VA: NCTM.

National Council of Teachers of Mathematics (NCTM) (2000).
\emph{Principles and Standards for School Mathematics}.
Reston, VA: NCTM.

National Council of Teachers of Mathematics (NCTM) (2006).
\emph{Curriculum Focal Points for Prekindergarten through Grade 8
Mathematics: A Quest for Coherence}.
Reston, VA: NCTM.

National Council of Teachers of Mathematics (NCTM) (2009).
\emph{Focus in High School Mathematics: Reasoning and Sense Making}.
Reston, VA: NCTM.

National Council of Teachers of Mathematics (NCTM) (2024).
\emph{Artificial Intelligence and Mathematics Teaching}.
Position Statement, February 2024. Reston, VA: NCTM.

Nguyen, D. T., và Tran, D. (2023).
``High school mathematics teachers' changes in beliefs and knowledge during
lesson study.''
\emph{Journal of Mathematics Teacher Education}, 26(6), 809--834.
doi: 10.1007/s10857-022-09547-2.

Trần Cường, Nguyễn Hoàng Anh, và Nguyễn Thị Hoài Thu (2024).
``Từ phong trào Toán học mới tới giáo dục toán học ở Việt Nam.''
\emph{Tạp chí Khoa học Trường Đại học Sư phạm Hà Nội: Khoa học Giáo dục},
69(5), 117--128.
doi: 10.18173/2354-1075.2024-0123.

Volkov, A. (2009).
``Mathematics and Mathematics Education in Traditional Vietnam.''
Trong E. Robson và J. Stedall (chủ biên),
\emph{The Oxford Handbook of the History of Mathematics},
tr. 153--176. Oxford University Press.
doi: 10.1093/oso/9780199213122.003.0008.

\endgroup
